\section{Detailed Software Design}
\label{sec:software-design}

\subsection{The layer stack (shared by both options)}

Twelve layers, listed bottom-up. Options A and B differ in \emph{how layers 1--3 get populated}; layers 4 upward are identical code paths regardless of origin.

\begin{enumerate}
\item \textbf{Runtime} --- the process(es) executing on the ACC (and, in a thin form, on each E2B node) that host drivers, scheduling, and the network stack. On the ACC: a Linux-hosted supervisor process plus a set of containerized "skills"; on RT domain and E2B nodes: a small deterministic executive (Zephyr-class or bare-metal) running generated driver code.
\item \textbf{Hardware Abstraction Layer (HAL)} --- per-interface-type drivers (PWM, ADC, SPI, CAN-FD, LIN, SENT, Ethernet) with a uniform internal API (\code{hal\_pwm\_set\_duty()}, \code{hal\_can\_send\_frame()}, etc.) so that the Verified Primitive Library above never touches a register directly.
\item \textbf{Verified Primitive Library} --- a curated set of driver \emph{instances} for specific hardware (a specific motor driver IC, a specific Hall sensor part, a specific CAN transceiver's init sequence) that have been through the verification pipeline in Section~\ref{sec:firmware-generation}. This is the layer that actually grows as a moat (Section~\ref{sec:moat}): each new peripheral supported adds one verified entry here, reusable on every future machine.
\item \textbf{Machine IR} (Section~\ref{sec:machine-ir}) --- the structured, versioned model of a specific machine instance: which device instances exist, which network path/resource binding each one uses, and (for legacy devices) which mapping confidence applies.
\item \textbf{Hardware Knowledge Graph} --- the cross-machine, cross-program database of \emph{general} facts: "this MCU part has N timer channels usable for PWM," "this transceiver part requires this init sequence," "this CAN ID pattern was previously seen meaning turn-signal on a different vehicle from the same OEM platform." The Machine IR for one vehicle references entries in this graph; the graph itself is shared across every vehicle Alfred has ever touched.
\item \textbf{Hardware Mapping / Device Abstraction} --- the layer that resolves a semantic capability call to a specific Machine-IR device instance and its bound primitive(s). This is literally the "Alfred hardware mapping" box in the architecture diagram in the background brief.
\item \textbf{Alfred Semantic API} --- the stable, typed verb surface (\code{door.open()}, \code{actuator.set\_position()}) that application code is written against.
\item \textbf{Behavior / Application Layer} --- the actual vehicle behaviors and business logic (e.g., "when door opens, if vehicle speed $>$ 0, ignore"), composed from Semantic API calls. This is where OEM/customer-specific logic lives, deliberately kept separate from everything below it.
\item \textbf{Firmware Compiler} (Section~\ref{sec:firmware-generation}) --- the deterministic toolchain that turns a Machine IR + placement decisions into buildable, flashable firmware images for the ACC RT domain and every E2B node.
\item \textbf{Target Packs} --- per-MCU/per-PHY "backend" definitions for the Firmware Compiler (pin maps, timer/DMA resource tables, linker scripts, board support package) --- adding a new MCU family to Alfred means writing one target pack, not touching the compiler.
\item \textbf{Deployment Manager} (Section~\ref{sec:deployment-reprogramming}) --- orchestrates SIL/HIL verification gates, staged rollout, and the actual programming of ACC and E2B targets, with rollback.
\item \textbf{Diagnostics \& Certification Evidence System} --- structured logging of every decision the pipeline made (resource allocation choices, test results, static analysis reports) in a form that is directly usable as functional-safety work-product evidence (ISO~26262 Part 6/8 traceability), not just human-readable debug output.
\end{enumerate}

\begin{diagram}
   Behavior / Application Layer  <-- OEM/customer logic; unchanged across A/B
              |
      Alfred Semantic API        <-- stable verbs
              |
    Hardware Mapping / Device Abstraction
              |
        +-----+------------------------+
        |                               |
   Machine IR (per-vehicle)     Hardware Knowledge Graph (cross-vehicle)
        |                               |
        +-------------+-----------------+
                      |
         Verified Primitive Library  <-- grows every program, reused forever
                      |
                     HAL
                      |
                   Runtime
\end{diagram}

\subsection{Option A: populating layers 1--4 directly}

An E2B node's on-boot manifest (Section~\ref{sec:discovery-a}) is generated \emph{by construction} --- the Firmware Compiler already knows, at build time, exactly which HAL drivers and Verified Primitive Library entries were compiled into that node's image, because it generated that firmware. The manifest the node reports at boot is therefore a confirmation/attestation of a known-good build, not a discovery in the Option-B sense. Machine IR population for a greenfield vehicle is consequently near-instant and 100\%-confidence: it is compiled alongside the firmware, not inferred after the fact.

\subsection{Option B: populating layers 1--4 through inference}

The Discovery software stack is the Option-B-specific software that does not exist in Option A, sitting logically \emph{below} the Machine IR layer as an alternative population path:

\begin{itemize}
\item \textbf{Capture Service} (per probe node, Section~\ref{sec:probe-discovery}) --- protocol-aware frame/event capture with hardware timestamping, writing to a common time-synchronized event log format.
\item \textbf{Protocol Identification Engine} --- Level 1/2 classification (Section~\ref{sec:probe-discovery}): electrical characterization, bitrate estimation, protocol decode (CAN$\rightarrow$ISO-TP$\rightarrow$UDS, Ethernet$\rightarrow$IP$\rightarrow$SOME/IP), implemented as deterministic decoders, not ML classifiers, because these are well-specified protocols with unambiguous grammars.
\item \textbf{Correlation Engine} --- Level 3 semantic identification: the statistical machinery (Section~\ref{sec:probe-discovery}) that turns "operator pressed the turn signal at T" plus "CAN-B ID 0x311 bit 4 changed at T+2ms" into a scored candidate mapping.
\item \textbf{Engineering File Importer} --- parsers for DBC, ARXML, LDF, A2L, ODX that pre-populate candidate signal names/scaling/units, which the Correlation Engine then uses as priors rather than starting from zero.
\item \textbf{Vehicle Knowledge Graph builder} (Section~\ref{sec:knowledge-graph}) --- assembles Capture Service + Protocol ID + Correlation Engine + imported files into the graph structure, which is what actually gets compiled into Machine IR entries once mappings cross a confidence threshold.
\item \textbf{Compiler: Graph $\rightarrow$ Machine IR} --- a deterministic step (not inference) that takes a knowledge-graph subtree whose confidence exceeds a configured threshold and mechanically emits a Machine IR device entry plus its network-path/hardware-mapping binding through a legacy protocol adapter (Section~\ref{sec:deployment-reprogramming}).
\end{itemize}

The key architectural point: everything above the Machine IR layer in Section~6.1 is identical code whether the IR entry was compiled directly (Option A) or produced by the Graph-to-IR compiler (Option B). The Correlation Engine and Graph builder are the only genuinely new software surfaces Option B requires beyond what Option A already needs.
